\subsection{Integration against a H\"older reference}

A H\"older modulus for a bounded potential gives a modulus for integration against its Monge--Amp\`ere measure. The proof uses a smooth approximation of the reference, with integration by parts controlling the change in its measure.

\begin{lemma}\label{lem:local-holder-reference-integrals}\index{Holder continuity@H\"older continuity!reference integrals}
Let $B'\Subset B$ be concentric balls in $\C^n$, and let $0<\alpha\leq1$. Suppose $p\in\PSH(B)\cap C^{0,\alpha}(B)$ has $C^{0,\alpha}$ norm at most $N$. For every $M\geq1$ there is a constant $C_M$, depending only on $M$, $N$, $\alpha$ and the balls, such that all $a,b\in\PSH(B)$ with $|a|,|b|\leq M$ satisfy
\[
 \int_{B'}|a-b|(\ddc p)^n
 \leq C_M\|a-b\|_{L^1(B)}^{\alpha/(\alpha+2n)}.
\]
\end{lemma}
\begin{proof}
Choose intermediate concentric balls and a smooth cutoff $0\leq\chi\leq1$ supported in the innermost of them and equal to one on $B'$. All integrations below take place away from $\partial B$. Let $p_\varepsilon$ be the convolution of $p$ with a nonnegative radial mollifier of radius $\varepsilon$, defined on a larger intermediate ball. It is plurisubharmonic, and, for small $\varepsilon>0$,
\[
 \|p-p_\varepsilon\|_\infty\leq C\varepsilon^\alpha,
 \qquad 0\leq\ddc p_\varepsilon\leq C\varepsilon^{-2}\ddc|z|^2.
\]
Both $p$ and $p_\varepsilon$ take values in $[-N,N]$. Set $f=|a-b|=2\max(a,b)-a-b$, and put
\[
 T=\sum_{j=0}^{n-1}(\ddc p)^j\wedge(\ddc p_\varepsilon)^{n-1-j}.
\]
This is a closed positive current. Factorization and integration by parts give
\[
 \int\chi f\bigl((\ddc p)^n-(\ddc p_\varepsilon)^n\bigr)
 =\int(p-p_\varepsilon)\ddc(\chi f)\wedge T.
\]
We check that the signed measure $\ddc(\chi f)\wedge T$ has total variation at most $C(1+M)^2$, uniformly in $\varepsilon$, where $C$ depends only on $N$, $\alpha$ and the balls. Expand
\[
 \ddc(\chi f)=\chi\ddc f+f\ddc\chi+\mathrm d\chi\wedge\dc f+\mathrm df\wedge\dc\chi.
\]
The term $\chi\ddc f\wedge T$ is a linear combination of the positive measures formed with $\max(a,b),a,b$, whose masses are at most $CM$ by local Chern--Levine--Nirenberg estimates. The term $f\ddc\chi\wedge T$ has variation at most $CM$.

For the two terms involving $\mathrm df$, use a cutoff $\eta$ equal to one near $\operatorname{supp}\chi$. If $q$ is any of $a,b,\max(a,b)$, translate it so that $0\leq q\leq2M$. The bounded chain rule and closedness of $T$ give
\[
 \int\eta\,\mathrm dq\wedge\dc q\wedge T
 =\frac12\int q^2\ddc\eta\wedge T
       -\int\eta q\ddc q\wedge T\leq CM^2.
\]
The last term is nonnegative before subtraction, and the first is bounded using the local mass of $T$. Since $f=2\max(a,b)-a-b$, its gradient energy is also bounded by $CM^2$. Cauchy--Schwarz against the positive current $T$ now bounds the variation of $\mathrm d\chi\wedge\dc f\wedge T$ and $\mathrm df\wedge\dc\chi\wedge T$ by $CM$. This accounts for all four terms of the expansion.

The bounded formulas follow first for local smooth regularizations. To pass to bounded potentials, regularize $a,b,p$ decreasingly on the intermediate balls, write $f$ as the displayed linear combination, and apply \cref{lem:weighted-decreasing-products} to the expanded identities. The energy bounds are uniform; thus the integration by parts and its estimate survive this limit.

Writing $L=\|a-b\|_{L^1(B)}$ and using $(\ddc p_\varepsilon)^n\leq C\varepsilon^{-2n}(\ddc|z|^2)^n$ for the second term, we have proved
\begin{equation}\label{eq:holder-reference-local-balance}
 \int\chi f(\ddc p)^n
 \leq C(1+M)^2\varepsilon^\alpha+C\varepsilon^{-2n}L.
\end{equation}
For small positive $L$, take $\varepsilon=L^{1/(\alpha+2n)}$. A fixed admissible radius and the bounded local mass handle the remaining $L$. For $L=0$, let $\varepsilon\downarrow0$ in the same inequality.
\end{proof}

\begin{corollary}\label{cor:holder-reference-weighted-modulus}
Let $(X,\omega)$ be compact Hermitian, and let $u\in\PSH(X,\omega)\cap C^{0,\alpha}(X)$, $0<\alpha\leq1$. If uniformly bounded quasi-plurisubharmonic functions $a,b$ have a fixed common smooth lower curvature bound, then
\[
 \int_X|a-b|\,\omega_u^n
 \leq C\|a-b\|_{L^1(\mathrm dV)}^{\alpha/(\alpha+2n)}.
\]
The constant depends on the bounds for $a,b$, their lower curvature bound and the $C^{0,\alpha}$ norm of $u$.
\end{corollary}
\begin{proof}
On each of finitely many coordinate balls, choose a fixed smooth strictly plurisubharmonic quadratic $Q$ such that $a+Q,b+Q,u+Q$ are plurisubharmonic and $\ddc Q\geq\omega$. There
$0\leq\omega_u\leq\ddc(u+Q)$, so positivity of bounded products gives $\omega_u^n\leq(\ddc(u+Q))^n$. Apply \cref{lem:local-holder-reference-integrals} on smaller balls covering $X$. The quadratic cancels in the difference, and the local volumes are equivalent to $\mathrm dV$. Summing proves the assertion.
\end{proof}

The estimate before optimization also controls the tails of unbounded potentials. We need the following local form with a compact ambient space, so that the exponential bound for smooth volume is uniform.

\begin{lemma}\label{lem:local-holder-exponential-tails}\index{Sublevel estimate!Holder reference@H\"older reference}
Let $(Y,\Omega)$ be compact Hermitian, let $B'\Subset B\subset Y$ be coordinate balls, and let $p\in\PSH(B)\cap C^{0,\alpha}(B)$ have norm at most $N$, with $0<\alpha\leq1$. There are $c,C>0$, uniform for these references, such that
\[
 (\ddc p)^n\bigl(B'\cap\{v<-t\}\bigr)\leq C\mathrm e^{-ct}
 \quad(t\geq0)
\]
for every $v\in\PSH(Y,\Omega)$ with $\sup_Yv=0$.
\end{lemma}
\begin{proof}
Choose a smooth quadratic $Q$ with $\ddc Q\geq\Omega$ on $B$ and a fixed constant $D$ such that $w=v+Q-D\leq0$. For integers $m\geq1$, set
\[
 a_m=\max(w,-m),\qquad b_m=\max(w,-m-1),\qquad f_m=a_m-b_m.
\]
These are bounded plurisubharmonic truncations, and $0\leq f_m\leq1$; moreover, $f_m$ vanishes on $\{w\geq-m\}$ and equals one on $\{w<-m-1\}$. By \cref{lem:uniform-exponential} on $(Y,\Omega)$ and the bounded correction $Q-D$,
\[
 \|f_m\|_{L^1(B)}\leq C\mathrm e^{-am}
\]
for a fixed $a>0$. Use the proof of \cref{lem:local-holder-reference-integrals}, with $M=m+1$, keeping its estimate \eqref{eq:holder-reference-local-balance} before optimization. Taking $\varepsilon=\mathrm e^{-am/(\alpha+2n)}$ gives
\[
 \int\chi f_m(\ddc p)^n
 \leq C(1+m)^2\exp\left(-\frac{a\alpha m}{\alpha+2n}\right)
 \leq C\mathrm e^{-cm}
\]
for some $c>0$. For large $m$ the convolution radius $\varepsilon$ is admissible, and the common local mass handles the other integers. Finally, $w-v=Q-D$ is bounded on $B$. Hence, for large $t$, the set $B'\cap\{v<-t\}$ lies in $\{w<-m-1\}$, where $\chi f_m=1$, for an integer $m\geq t-C$. This proves the estimate for large real $t$, and the common local mass covers the remaining $t\geq0$.
\end{proof}
